\documentclass[aspectratio=169,10pt]{ctexbeamer}

% 建议使用 XeLaTeX 或 LuaLaTeX 编译。
\usepackage{amsmath,amssymb}
\usepackage{booktabs}
\usepackage{tikz}
\usepackage{hyperref}
\usetikzlibrary{positioning,arrows.meta}

% 中文字体：正文使用宋体。建议使用 XeLaTeX 或 LuaLaTeX 编译。
\setCJKmainfont{SimSun}
\setCJKsansfont{SimSun}
\usefonttheme{serif}

% ---------- 清爽的自定义样式 ----------
\definecolor{DeepBlue}{HTML}{183247}
\definecolor{MutedBlue}{HTML}{2D6280}
\definecolor{SoftTeal}{HTML}{478C88}
\definecolor{WarmRed}{HTML}{B45C3E}
\definecolor{WarmGold}{HTML}{C49A4A}
\definecolor{Ink}{HTML}{263238}
\definecolor{Muted}{HTML}{607078}
\definecolor{Pale}{HTML}{EDF3F3}

\setbeamercolor{normal text}{fg=Ink,bg=white}
\setbeamercolor{structure}{fg=SoftTeal}
\setbeamercolor{background canvas}{bg=white}
\setbeamercolor{frametitle}{fg=DeepBlue,bg=white}
\setbeamercolor{title}{fg=DeepBlue}
\setbeamercolor{subtitle}{fg=MutedBlue}
\setbeamercolor{alerted text}{fg=WarmRed}
\setbeamercolor{block title}{fg=DeepBlue,bg=Pale}
\setbeamercolor{block body}{fg=Ink,bg=Pale}

\setbeamertemplate{navigation symbols}{}
\setbeamertemplate{blocks}[rounded][shadow=false]
\setbeamertemplate{itemize item}{\color{SoftTeal}\small$\blacktriangleright$}
\setbeamertemplate{itemize subitem}{\color{WarmGold}\scriptsize$\bullet$}
\setbeamertemplate{frametitle}{%
  \vspace{0.30cm}\hspace{0.15cm}\insertframetitle\par
  \vspace{0.10cm}\hspace{0.15cm}\color{SoftTeal}\rule{1.35cm}{1.5pt}
}
\setbeamertemplate{footline}{%
  \leavevmode
  \hbox{%
    \begin{beamercolorbox}[wd=.86\paperwidth,ht=2.7ex,dp=1.2ex,leftskip=.55cm]{author in head/foot}%
      \color{Muted}\scriptsize 自然辩证法课程汇报
    \end{beamercolorbox}%
    \begin{beamercolorbox}[wd=.14\paperwidth,ht=2.7ex,dp=1.2ex,rightskip=.55cm plus1fil]{page number in head/foot}%
      \color{SoftTeal}\scriptsize\insertframenumber/\inserttotalframenumber
    \end{beamercolorbox}%
  }
}

\setlength{\leftmargini}{1.3em}
\setlength{\leftmarginii}{1.1em}
\newcommand{\key}[1]{\textcolor{DeepBlue}{\textbf{#1}}}
\newcommand{\warn}[1]{\textcolor{WarmRed}{\textbf{#1}}}
\newcommand{\source}[1]{%
  \par\vfill\begingroup\color{Muted}\fontsize{5.8}{7}\selectfont 来源：#1\endgroup
}
\newcommand{\takeaway}[1]{%
  \begin{beamercolorbox}[rounded=true,sep=0.65em,wd=\linewidth]{block body}
    \centering\color{MutedBlue}\textbf{#1}
  \end{beamercolorbox}
}

\title{科研越轨的成因与治理}
\subtitle{基于跨学科案例的认识论与制度分析}
\author{汇报人：\underline{\hspace{3.2cm}}}
\institute{自然辩证法课程汇报}
\date{学号：\underline{\hspace{3.2cm}}}

\begin{document}

% 1
\begin{frame}[plain]
  \titlepage
\end{frame}

% 2
\begin{frame}{科研越轨的边界}
  \begin{columns}[T,onlytextwidth]
    \column{0.31\textwidth}
      \begin{block}{诚实错误}
        方法局限、随机波动或无意失误。研究者应当及时更正，并保留可追溯记录。
      \end{block}
    \column{0.31\textwidth}
      \begin{block}{有害研究实践}
        选择性报告、数据管理不善、署名不当等。未必达到失信认定标准，却会削弱知识可靠性。
      \end{block}
    \column{0.31\textwidth}
      \begin{block}{科研失信}
        伪造、篡改、剽窃，以及买卖论文、虚假评审、无实质贡献署名等明确违规行为。
      \end{block}
  \end{columns}

  \vspace{0.45cm}
  \takeaway{判断不能只看结论是否错误，还要考察研究过程、行为责任和科研记录是否受到破坏。}
  \source{美国研究诚信办公室 ORI；科技部等《科研失信行为调查处理规则》（2022）}
\end{frame}

% 3
\begin{frame}{“天使粒子”论文的兴衰}
  \begin{tikzpicture}[remember picture,overlay]
    \draw[SoftTeal,line width=1.2pt] (1.65,-0.25) -- (1.65,-5.45);
    \foreach \y/\c in {-0.55/SoftTeal,-2.00/WarmGold,-3.45/WarmRed,-4.90/WarmRed}
      \fill[\c] (1.65,\y) circle (3.1pt);
  \end{tikzpicture}

  \begin{description}[2022.11]
    \item[2017.07] 《Science》论文报告接近 $0.5e^2/h$ 的电导平台，并将其解释为手性马约拉纳模式。
    \vspace{0.24cm}
    \item[2020.01] 两个独立实验团队未观察到相应证据，并指出原实验几何可能造成电流短路。
    \vspace{0.24cm}
    \item[2021.12] 《Science》发布编辑关注声明，提醒读者数据与结论仍存在疑问。
    \vspace{0.24cm}
    \item[2022.11] 编辑指出原始数据来源受到质疑，原始数据与发表数据存在严重异常和不一致，决定撤稿。
  \end{description}

  \source{He et al., Science 357, 294--299 (2017)；Science Editorial Retraction, DOI: 10.1126/science.adf7575}
\end{frame}

% 4
\begin{frame}{“天使粒子”案例的证据争议}
  \begin{columns}[T,onlytextwidth]
    \column{0.31\textwidth}
      \key{论文中的关键证据}
      \vspace{0.15cm}

      磁场扫描中出现接近半量子化的电导平台。作者将平台与磁化反转联系起来，视为手性马约拉纳模式的特征信号。
    \column{0.31\textwidth}
      \key{后续检验的问题}
      \vspace{0.15cm}

      独立实验未能复现，并提出短路等替代解释。同行取得数据后，又发现数据来源和处理过程存在严重疑点。
    \column{0.31\textwidth}
      \key{证据意义的变化}
      \vspace{0.15cm}

      半量子化平台无法排除其他机制。原始数据不能稳定追溯时，同行也无法判断公开图表是否忠实代表实验记录。
  \end{columns}

  \vspace{0.45cm}
  \takeaway{复现失败本身不能证明科研不端；数据来源不清和处理不可追溯，才使问题进入科研诚信审查。}
  \source{Chang et al., Science 367, 64--67 (2020)；Science Editorial Retraction (2022)}
\end{frame}

% 5
\begin{frame}{“天使粒子”案例的成因}
  \begin{columns}[T,onlytextwidth]
    \column{0.48\textwidth}
      \begin{block}{认识论层面}
        \begin{itemize}
          \item 研究者依据电导平台间接推断准粒子模式，相同信号可能对应多种物理机制。
          \item 当 $0.5e^2/h$ 成为预期标志后，符合理论的扫描区间更容易受到重视。
          \item 局部“漂亮数据”可能掩盖完整参数范围中的反例。
        \end{itemize}
      \end{block}
    \column{0.48\textwidth}
      \begin{block}{制度层面}
        \begin{itemize}
          \item 拓扑量子计算的战略价值放大了首发成果的声誉与资源收益。
          \item 多团队合作将数据生产、处理、理论解释和署名责任分散开来。
          \item 期刊与媒体偏好明确的突破叙事，弱化了不确定性表达。
        \end{itemize}
      \end{block}
  \end{columns}

  \vspace{0.35cm}
  \takeaway{公布完整扫描范围和处理脚本，并明确数据保管人与共同作者的核验责任。}
\end{frame}

% 6
\begin{frame}{舍恩事件：重复曲线如何暴露造假}
  \begin{enumerate}
    \item \key{突破性成果密集出现}：短期内连续发表有机半导体与分子器件成果，多个结果呈现近乎理想的曲线。
    \item \key{异常来自跨论文比较}：研究者发现不同论文、不同实验条件下的噪声形态和数据曲线异常相似。
    \item \key{调查回到原始记录}：贝尔实验室外部委员会核查原始数据、图表制作过程和共同作者责任。
    \item \key{调查结论}：委员会在所审查事项中认定多项科研不端，贝尔实验室随后解除舍恩的职务。
  \end{enumerate}

  \vspace{0.30cm}
  \begin{columns}[T,onlytextwidth]
    \column{0.48\textwidth}
      \warn{认识论教训}：同行评议主要检查论证是否合理，无法替代原始数据审计。
    \column{0.48\textwidth}
      \warn{制度教训}：成果收益集中于个人，数据保存和共同作者核验责任却容易被稀释。
  \end{columns}
  \source{Lucent Technologies, Report of the Investigation Committee on the Work of Hendrik Schön and Coauthors (2002)}
\end{frame}

% 7
\begin{frame}{STAP 细胞论文：从图像异常到样本失真}
  \begin{description}[后续验证]
    \item[2014 年 1 月] 两篇论文声称外界刺激可使体细胞获得多能性。
    \item[质疑出现] 同行发现图像、文字和实验复现方面的异常。
    \item[调查撤稿] 理化学研究所开展调查，两篇论文于 2014 年撤回。
    \item[后续验证] 基因组分析显示所谓 STAP 细胞系来自胚胎干细胞污染；独立实验也未能复现该现象。
  \end{description}

  \vspace{0.35cm}
  \takeaway{图片检查只能发现表征层异常，样本谱系、实验记录和独立复现才能检验核心主张。}
  \source{RIKEN investigation reports (2014)；Konno et al., Nature 525, E4--E5 (2015)；Aizawa et al., Nature 525, E6--E9 (2015)}
\end{frame}

% 8
\begin{frame}{认识论成因：证据链如何断裂}
  \begin{center}
  \begin{tikzpicture}[
    node distance=0.42cm,
    stage/.style={draw=SoftTeal!55,fill=Pale,rounded corners=2pt,
                  minimum width=2.25cm,minimum height=1.05cm,align=center},
    arr/.style={-{Latex[length=2mm]},SoftTeal!75,line width=.8pt}]
    \node[stage] (a) {问题与假设\\\scriptsize 追逐热点、确认偏误};
    \node[stage,right=of a] (b) {采集与记录\\\scriptsize 原始数据缺失};
    \node[stage,right=of b] (c) {处理与统计\\\scriptsize 选择区间、过度拟合};
    \node[stage,right=of c] (d) {解释与发表\\\scriptsize 叙事先行};
    \node[stage,right=of d] (e) {重复与批判\\\scriptsize 负结果不传播};
    \draw[arr] (a)--(b); \draw[arr] (b)--(c); \draw[arr] (c)--(d); \draw[arr] (d)--(e);
  \end{tikzpicture}
  \end{center}

  \vspace{0.35cm}
  \begin{itemize}
    \item \key{确认偏误}使研究者偏爱支持假设的证据。
    \item \key{权威效应}使声望替代对原始证据的审查。
    \item \key{结果导向}以显著结论评价研究过程，压低不确定性和失败实验的价值。
  \end{itemize}
  \source{National Academies, Fostering Integrity in Research (2017)}
\end{frame}

% 9
\begin{frame}{制度成因：激励、监督与权力结构}
  \small
  \begin{tabular}{@{}p{0.15\linewidth}p{0.42\linewidth}p{0.36\linewidth}@{}}
    \toprule
    \textbf{制度环节} & \textbf{具体表现} & \textbf{行为后果} \\
    \midrule
    评价激励 & 论文数量、期刊层级与短周期考核过度绑定 & 显著结果和快速成功获得更高回报 \\
    \addlinespace
    组织监督 & 原始数据无人复核，共同作者责任边界模糊 & 违规被发现的概率和成本降低 \\
    \addlinespace
    资源竞争 & 职位、经费和项目稀缺，失败容忍度低 & 选择性报告与抢发倾向增强 \\
    \addlinespace
    权力关系 & 学生依附导师，举报渠道缺少独立性 & 知情者沉默，问题长期积累 \\
    \addlinespace
    出版机制 & 负结果难发表，期刊偏好新奇和显著性 & 文献形成系统性偏差 \\
    \bottomrule
  \end{tabular}
  \source{National Academies, Fostering Integrity in Research (2017)}
\end{frame}

% 10
\begin{frame}{认识因素与制度因素相互放大}
  \begin{center}
  \begin{tikzpicture}[
    node distance=0.9cm and 1.35cm,
    cause/.style={draw=SoftTeal!60,rounded corners=2pt,fill=white,
                  minimum width=3.2cm,minimum height=.85cm,align=center},
    risk/.style={fill=DeepBlue,text=white,rounded corners=3pt,
                 minimum width=3.7cm,minimum height=1.2cm,align=center},
    arr/.style={-{Latex[length=2mm]},line width=1pt}]
    \node[risk] (r) {\large 越轨风险上升\\\scriptsize 错误被掩盖，违规获得收益};
    \node[cause,above left=of r] (a) {证据不透明\\\scriptsize 原始数据不可追溯};
    \node[cause,below left=of r] (b) {确认偏误\\\scriptsize 异常结果缺少反证};
    \node[cause,above right=of r] (c) {评价单一\\\scriptsize 数量和期刊指标主导};
    \node[cause,below right=of r] (d) {权力不对称\\\scriptsize 质疑与举报成本高};
    \draw[arr,SoftTeal] (a)--(r); \draw[arr,WarmGold] (b)--(r);
    \draw[arr,WarmRed] (c)--(r); \draw[arr,WarmRed] (d)--(r);
  \end{tikzpicture}
  \end{center}

  \vspace{0.3cm}
  \takeaway{有效治理必须同时修复知识生产过程和组织激励结构。}
\end{frame}

% 11
\begin{frame}{认识论治理：缩短错误暴露时间}
  \begin{enumerate}
    \item \key{全过程留痕}：保存原始数据、样本来源、分析代码和版本记录。
    \item \key{预先约束分析自由度}：预注册主要假设、指标和排除规则，区分探索性分析与验证性分析。
    \item \key{关键结果独立复核}：由未直接参与分析的研究者复算关键图表，核查样本和统计流程。
    \item \key{提高可重复性}：共享数据、代码、材料和实验协议，支持重复实验和负结果传播。
    \item \key{常态化批判训练}：将统计推断、科研记录、图像规范和案例讨论纳入研究生培养。
  \end{enumerate}
  \source{UNESCO, Recommendation on Open Science (2021)；National Academies (2017)}
\end{frame}

% 12
\begin{frame}{制度治理：重塑激励并建立独立问责}
  \begin{columns}[T,onlytextwidth]
    \column{0.48\textwidth}
      \begin{block}{评价与组织责任}
        \begin{itemize}
          \item 减少对论文数量和期刊层级的机械依赖。
          \item 评价数据质量、方法贡献和纠错记录。
          \item 明确导师、共同作者和数据保管人的责任。
          \item 对高风险研究环节开展内部抽查。
        \end{itemize}
      \end{block}
    \column{0.48\textwidth}
      \begin{block}{调查程序与外部协同}
        \begin{itemize}
          \item 使受理、调查与行政利益保持相对独立。
          \item 保护善意举报者，也保障被举报者的程序权利。
          \item 基金方、机构和期刊共享必要线索。
          \item 及时更正或撤回论文，修复公开科研记录。
        \end{itemize}
      \end{block}
  \end{columns}

  \vspace{0.25cm}
  \takeaway{制度应当让守规矩者不吃亏，让提出质疑者不受惩罚，让违规收益低于预期成本。}
  \source{科技部等《科研失信行为调查处理规则》（2022）；National Academies (2017)}
\end{frame}

% 13
\begin{frame}{跨学科科研的操作清单}
  \small
  \begin{tabular}{@{}p{0.18\linewidth}p{0.76\linewidth}@{}}
    \toprule
    \textbf{研究环节} & \textbf{最低要求} \\
    \midrule
    实验与样本 & 原始文件只读归档，记录仪器校准、样本来源、批次信息和失败实验。\\
    \addlinespace
    图像与组学 & 保存未处理数据，记录裁剪、归一化、背景扣除和生物信息流程。\\
    \addlinespace
    模型与计算 & 锁定代码版本、依赖环境、训练数据、随机种子和输入参数。\\
    \addlinespace
    临床与调查 & 预注册主要终点，保存知情同意，说明脱落样本和缺失数据处理。\\
    \addlinespace
    论文与署名 & 逐项确认贡献与责任，所有作者核读关键图表和主要结论。\\
    \addlinespace
    独立复核 & 投稿前安排交叉复算、复测或盲审，高影响结论优先开展外部复现。\\
    \bottomrule
  \end{tabular}
\end{frame}

% 14
\begin{frame}{治理行动的优先顺序}
  \small
  \begin{tabular}{@{}p{0.11\linewidth}p{0.16\linewidth}p{0.40\linewidth}p{0.25\linewidth}@{}}
    \toprule
    \textbf{时间} & \textbf{主体} & \textbf{核心行动} & \textbf{预期作用} \\
    \midrule
    近期 & 实验室 & 原始数据规范、投稿前复核、署名确认 & 降低无意错误和低成本违规 \\
    \addlinespace
    中期 & 院系与机构 & 独立举报渠道、抽查审计、导师责任 & 提高发现概率，降低沉默成本 \\
    \addlinespace
    中期 & 评价与资助部门 & 代表作评价、过程质量、纠错记录 & 削弱只奖励结果的单一激励 \\
    \addlinespace
    长期 & 科研共同体 & 数据标准、期刊联动、复现基础设施 & 恢复文献系统的自我纠错能力 \\
    \bottomrule
  \end{tabular}

  \vspace{0.45cm}
  \takeaway{认识论治理守住证据底线，评价改革改变行为激励，独立调查提供最后保障。}
\end{frame}

\end{document}
